
\subsubsection{6.1 The Gap Between Mechanism and Manifestation}

The mechanistic differences between RLHF and DPO were verified (H-M1, H-M2), but these did not produce distinct behavioral signatures (H-M3, H-M4). Several factors may explain this gap:

\textbf{Seed variance dominance:} At 7B scale with the evaluated training configuration, random seed variation appears to dominate over training method variation. Models trained with the same method but different seeds were less similar to each other than models trained with different methods.

\textbf{LoRA constraints:} The low-rank adaptation may constrain the parameter space sufficiently that both methods converge to similar regions despite different optimization paths. Full fine-tuning might permit larger divergence.

\textbf{Benchmark granularity:} Standard benchmarks may aggregate over dimensions in ways that mask method-specific differences. Item-level analysis or custom probes might reveal finer-grained distinctions.

\textbf{Convergence dominance:} The shared optimization objective (aligning with human preferences in HH-RLHF) may dominate over the different paths taken to reach it. Both methods ultimately optimize for the same preference signal.

\subsubsection{6.2 Implications for Method Selection}

At 7B scale with LoRA fine-tuning, the choice between RLHF and DPO may matter less for final model behavior than previously assumed. The mechanistic differences are real but appear to wash out by evaluation time. Practitioners may reasonably select based on computational convenience (DPO avoids training a separate reward model) rather than expected behavioral differences.

\subsubsection{6.3 Evaluation Sensitivity}

The high profile correlation (0.978) between DPO and RLHF on standard benchmarks suggests these benchmarks may not be sensitive to method-specific alignment differences. More granular evaluation approaches---item-level analysis, specialized probes, or custom metrics targeting known mechanistic differences---may be needed to detect such differences if they exist.

\subsubsection{6.4 Limitations}

\textbf{Quick validation mode:} H-M3 and H-M4 used abbreviated validation due to compute constraints. Results require replication with full multi-seed GPU training.

\textbf{Scale constraint:} All experiments used Llama-2-7B with LoRA. Larger models (13B, 70B) or full fine-tuning may show different patterns.

\textbf{Single dataset:} Only Anthropic HH-RLHF was used. Dataset characteristics may influence the comparative behavior of the methods.

\textbf{Single-epoch training:} Longer training might amplify method differences that are subtle after one epoch.

\textbf{Simulation mode:} H-M4 ran on simulated data because H-M3 did not produce trained checkpoints during execution. The benchmark comparison results require replication with actual trained models.

